\documentclass[11pt]{article}
\usepackage[margin=1in]{geometry}
\usepackage{amsmath,amssymb}
\usepackage{graphicx}
\usepackage{microtype}
\usepackage[hidelinks]{hyperref}
\setlength{\emergencystretch}{2em}

\newcommand{\Pd}{\mathbb P_d}
\newcommand{\Gatt}{\mathcal G_d}
\newcommand{\Uatt}{\mathcal U_{\mathrm{att}}}

\title{\textbf{The Global Set of Attainable Unitary Factors}\\
A Literature-Implied Complete Characterization}
\author{Status note for arXiv:2605.28885}
\date{12 August 2026}

\begin{document}
\maketitle

\begin{center}
\fbox{\parbox{0.91\textwidth}{
\textbf{Status: literature-implied answer.}
This is an agent-identified consequence of a 1986 matrix-factorization
theorem, not an original result and not an answer knowingly directed at the
2026 source paper.
}}
\end{center}

\section{Source question}

Vuong's arXiv:2605.28885 proves a fixed-pair criterion for the unitary factors
of generalized fidelity \cite{Vuong2026}.  In the Discussion (PDF p.~27), the
paper identifies
\[
 \Gatt:=
 \bigcup_{P,Q\in\Pd}
 \{W\in U(d):P^{-1/2}Q^{1/2}W\text{ is similar to a positive-definite
 matrix}\}
\tag{1}\label{eq:G}
\]
as the natural global object to characterize more explicitly.

\begin{figure}[ht]
\centering
\includegraphics[width=\textwidth]{figures/open_question_crop.png}
\caption{The stated global object in arXiv:2605.28885, p.~27.}
\end{figure}

\section{Identification with a classical factorization class}

We use the standard fact, also implicit in the source paper, that a complex
matrix is similar to a positive-definite matrix if and only if it is a product
of two positive-definite matrices.

\paragraph{Attainable implies four factors.}
If $W\in\Gatt$, put $M=P^{-1/2}Q^{1/2}$ for a witnessing pair.  Both $M$ and
$MW$ are products of two positive-definite matrices.  Since the inverse of a
positive-definite matrix is positive definite,
\[
 W=M^{-1}(MW)
\]
is a product of four positive-definite matrices.

\paragraph{Four factors imply attainable.}
Conversely, suppose
\[
 W=H_1H_2H_3H_4,\qquad H_j>0.
\]
Set
\[
 P=H_2^2,\qquad Q=H_1^{-2}.
\]
Then
\[
 M=P^{-1/2}Q^{1/2}=H_2^{-1}H_1^{-1},
\qquad
 MW=H_3H_4.
\]
The latter is a product of two positive-definite matrices and hence is similar
to a positive-definite matrix.  The fixed-pair criterion of
\cite[Theorem 5.9]{Vuong2026} therefore makes $W$ attainable.  We have proved
\[
 W\in\Gatt
 \quad\Longleftrightarrow\quad
 W\text{ is a product of four positive-definite matrices}.
\tag{2}\label{eq:four}
\]

\section{Sourour--Ballantine criterion and the answer}

Sourour's Theorem~2, attributed there to Ballantine, states that a real or
complex matrix $A$ is a product of four positive-definite matrices if and only
if
\[
 \det A>0
 \quad\text{and}\quad
 A\ne\alpha I\ \text{for every scalar $\alpha$ that is not positive}.
\tag{3}\label{eq:ballantine}
\]
Here $\det A>0$ means a positive real determinant \cite{Sourour1986}.

\begin{figure}[ht]
\centering
\includegraphics[width=\textwidth]{figures/supporting_theorem_crop.png}
\caption{Sourour \cite[Theorem 2]{Sourour1986}, reproducing Ballantine's
four-positive-factor criterion.}
\end{figure}

For a unitary $W$, the first condition in \eqref{eq:ballantine} is equivalent
to $\det W=1$, namely $W\in SU(d)$.  The scalar elements of $SU(d)$ are
$\omega I$ with $\omega^d=1$.  Among them, only $I$ has a positive scalar.
Combining \eqref{eq:four} and \eqref{eq:ballantine} gives the complete
characterization
\[
 \boxed{
 \Gatt
 =
 SU(d)\setminus
 \{\omega I:\omega^d=1,\ \omega\ne1\}.
 }
\tag{4}\label{eq:answer}
\]

This sharpens Corollary~5.15 of the source: the nonidentity central special
unitaries are not merely examples of excluded factors; they are the
\emph{only} excluded elements of $SU(d)$.

\section{Provenance and remaining scope}

The supporting theorem predates the source by forty years and does not mention
quantum fidelity.  The source paper does not cite Ballantine or Sourour.
Accordingly, \eqref{eq:answer} is classified as a
\emph{literature-implied answer}: the mathematical input is known, while the
identification with \eqref{eq:G} is the bridge recorded here.

The bounded search checked the run's cheap indexes; the source's references;
exact searches for products of four positive-definite matrices; Ballantine's
1968/1970 papers; Sourour's primary 1986 paper; and arXiv:2507.12560 (2025),
which surveys constructive positive-factor decompositions.  No evidence was
found that the 2026 source authors had already made this identification.

This note resolves only the explicit global-unitary-factor object.  Full
data-processing monotonicity, related convexity questions, and the proposed
semidefinite-program formulations remain open.

\begin{thebibliography}{9}

\bibitem{Vuong2026}
T.~D.~Vuong,
\emph{Polar Fidelities, Holevo Bases, and Unitary Factors of Generalized
Fidelity},
arXiv:2605.28885, 2026.

\bibitem{Sourour1986}
A.~R.~Sourour,
\emph{A Factorization Theorem for Matrices},
Linear and Multilinear Algebra \textbf{19}(2) (1986), 141--147,
\href{https://doi.org/10.1080/03081088608817711}{doi:10.1080/03081088608817711}.

\end{thebibliography}

\end{document}
